\section{Introduction}

Lending and borrowing protocols are among the most widely used blockchain applications, totaling about \$25 billion of active loans as of mid-2026~\cite{defillama}. Pool-based markets provided an effective design in early environments of thin liquidity and high transaction costs: all supply aggregated, rates set algorithmically from utilization, positions entered and exited at will~\cite{aave2022,compound2019}. The model works well for variable-rate, open-term loans---but cannot naturally express fixed-term or fixed-rate primitives, and commingles risk across all positions in the pool. As perpetual futures markets mature, the infrastructure supporting them is similarly evolving---from exchange-native margin systems where traders bring their own capital toward brokerage models that separate capital provision from trade execution.

A holder of a single collateral asset on one chain cannot easily obtain credit exposure across multiple venues and products without manually composing across protocols---bearing the bridging, timing, and counterparty risk of each step. No unified platform currently manages positions across chains, converts variable-rate exposure to fixed terms, or provides a programmatic interface for leveraged strategies that preserve the native yield of the pledged asset. Existing lending protocols handle individual positions within shared pools; protocols like Morpho Blue and Midnight provide isolated, permissionless markets---fixed-maturity in Midnight's case---but each operates on a single chain and a single market at a time~\cite{morphoblue2023,midnight2025}. Concurrent approaches to cross-venue prime brokerage aggregate all positions into a unified margin account, gaining capital efficiency from cross-venue position netting but coupling all venue exposure into a single solvency computation---a failure at any connected venue can cascade into global liquidation. None of these designs simultaneously aggregates liquidity across venues, converts variable-rate funding to fixed terms with hedged rate risk, and orchestrates cross-chain collateral and venue exposure on behalf of the user. Further, institutional counterparties with jurisdictional or compliance requirements cannot participate in permissionless order books without either fragmenting into permissioned forks---losing unified liquidity---or ignoring compliance entirely.

\daloc{} is a multichain lending platform with three user-facing products built on a common platform layer. \textbf{Borrow} provides fixed-rate, fixed-term loans with cross-chain collateral. \textbf{Prime Trade}---the platform's prime brokerage product on Hyperliquid---provides leveraged credit lines for perpetual trading while preserving the borrower's native yield. \textbf{Earn} lets liquidity providers supply capital into the Harvest Vault, where an operator deploys it into loans under a cryptographically-enforced mandate.

Underneath, the Harvest Vault is the platform's primary capital source: when it has sufficient idle capital, it funds loans directly at the full fixed rate, capturing the complete spread. When vault liquidity is fully deployed, external lending venues---Aave, Morpho, Spark, and others---provide overflow capacity through atomic hook chains, with the vault committing only hedging buffers that convert variable rates to fixed terms. This dual-path model gives the platform lending depth from its first deposit and scales without ceiling or dependency on token emissions.

An order book settles all loan activity---borrower requests and solver quotes are symmetric entries in the same book, matched and committed through epoch-based settlement verified by zero-knowledge proofs. The matching engine evaluates criteria-based compatibility as part of the match validity check, enabling restricted and unrestricted counterparties to share a single liquidity surface without permissioned forks. All external interactions within a loan are composed from atomic hook chains on the EVM, enabling multi-step pipelines---deposit, supply, borrow, bridge---within a single transaction. Markets are defined by pluggable zero-knowledge programs that can encode anything from classical locked-collateral rules to capital-efficient strategies permitting managed rehypothecation, with all logic publicly visible and cryptographically enforced.
